\clearpage
\subsection{Temperaturbestimmung -- Aufgabe 2}

Gegeben sind ein ADC-Messwert von $N=520$, eine Aufl\"osung von
$n=10\,\mathrm{Bit}$ und eine maximale ADC-Eingangsspannung von
$U_{\mathrm{ADC,max}}=5\,\mathrm{V}$. Der lineare Temperatursensor liefert bei
$T_{\max}=40\,^\circ\mathrm{C}$ eine Spannung von
$U_{\mathrm{Sensor,max}}=3{,}3\,\mathrm{V}$. Gesucht ist die zum Messwert
geh\"orende Temperatur $T$.

Zun\"achst wird der ADC-Messwert in eine Eingangsspannung umgerechnet. Ein
$n$-Bit-ADC besitzt $2^n$ Abstufungen; der gr\"o{\ss}te darstellbare Messwert
ist daher $N_{\max}=2^n-1$. Anschlie{\ss}end wird die Spannung anhand der
linearen Kennlinie des Sensors in eine Temperatur umgerechnet:

\begin{align*}
N_{\max}
    &= 2^n - 1
     = 2^{10} - 1
     = 1023, \\
U_{\mathrm{Sensor}}
    &= \frac{N}{N_{\max}}\,U_{\mathrm{ADC,max}} \\
    &= \frac{520}{1023}\cdot 5\,\mathrm{V}
     = 2{,}5415\,\mathrm{V}
     \approx 2{,}54\,\mathrm{V}, \\
T
    &= \frac{U_{\mathrm{Sensor}}}{U_{\mathrm{Sensor,max}}}\,T_{\max} \\
    &= \frac{2{,}5415\,\mathrm{V}}{3{,}3\,\mathrm{V}}
       \cdot 40\,^\circ\mathrm{C}
     = 30{,}806\ldots\,^\circ\mathrm{C}
     \approx 30{,}8\,^\circ\mathrm{C}.
\end{align*}

\noindent\textbf{Ergebnis:}
\[
    \boxed{T \approx 30{,}8\,^\circ\mathrm{C}}
\]
Der ADC-Messwert 520 entspricht somit einer Temperatur von etwa
$30{,}8\,^\circ\mathrm{C}$.
